\scanpage{82}
% archive sheet number: 81
% block text | lines 1-2
Quel est \unsure{donc} $M_f^\Gamma$ ? C'est l'ensemble des $m\in M_f=M$ tels que
% block display | line 3
\[
f(\unsure{\gamma|\sigma})\gamma m=m\qquad\text{pour }\unsure{\mathrm{tt}|\forall}\ \gamma\qquad\text{i.e.}
\]
% block display | line 4
\[
\boxed{\gamma m=f(\gamma)^{-1}m\qquad\text{pour }\forall\ \gamma\in\Gamma}
\]
% block text | lines 5-6
On les appelle parfois éléments $f$-automorphes du $G$-$\Gamma$-ensemble $M$.

% block text | lines 7-13
Si $M$ est muni \inserted{en plus} de structures diverses invariantes par $\Gamma$ et $G$, alors il en est de $=$ de $M_f$ (qui conserve les structures qu'il y avait sur $M$). Si $M$ était p.ex. un $G$-groupe abélien, \deleted{\unsure{il est|et}} un $\Gamma$-groupe abélien, il en est de $=$ de $M_f$ etc. Dans ce cas, on peut s'intéresser, plus généralement, aux $H^i(\Gamma,M_f)$.

% block text | lines 14-17
Soit $f'$ \unsure{une|un} autre \unsure{$1$-cocycle|cocycle} de $\Gamma$ à valeurs dans $G$. Supposons $f'\sim f$, i.e. $f'\gamma=a^{-1}f(\gamma)\gamma a\quad(a\in G)$. On a alors un isomorphisme $u_a:M_{f'}\longrightarrow M_{f\deleted{\unsure{'}}}$, en \deleted{faisant correspondre}\inserted{posant}
% block display | line 18
\[
u_a m=a m
\]
% block text | line 19
Il faut vérifier
% block display | lines 20-22
\[
\begin{aligned}
u_a\gamma_{f'}m&=\gamma_f u_a m\\
\text{i.e.}\qquad a f'(\gamma)\gamma m&=f^{\deleted{\unsure{'}}}(\gamma)\gamma(am)\\
\text{i.e.}\qquad aa^{-1}f(\gamma)\gamma(a)\gamma(m)&=f(\gamma)\gamma(am)\ ;\ \text{c'est vrai.}
\end{aligned}
\]
% block text | lines 23-25
Si on veut que $\unsure{M_g|M_f}\simeq M_{f'}$ pour $\mathrm{tt}$ $G$-$\Gamma$-ensemble $M$, il \uline{faut} que $f\sim f'$ comme on voit en prenant $M=G$, $G$ y opérant par \unsure{translations gauches}.
